\subsection{\texttt{TCOLEXPANDMIN}}
\noindent\textbf{Bundle encoding.} UOP entry 10, opcode \texttt{0x1015}. Bundle: \texttt{B.OP + B.ARG + B.DIM + B.IOTI}. \texttt{B.OP.Opcode}=\texttt{0x1015}. \texttt{B.IOTI}: selector=\texttt{111 last}, \texttt{SrcTile0}=\texttt{src0}, \texttt{SrcTile1}=\texttt{src1}, \texttt{DstTile}=\texttt{dst}, \texttt{SizeCode}=profile tile size. \texttt{B.DIM} supplies column-axis reduction or column-broadcast bounds.\par\smallskip
{\small
\paragraph{PTO ISA description.}
Column-wise broadcast min: take \texttt{min(src0, src1)} where \texttt{src1} provides one scalar per column.
\paragraph{Example.}
dst[:,j] = min(tile[:,j], src[0,j]) for each column j.
\par\smallskip
\paragraph{PTO ISA operation.}
Let \texttt{R = dst.GetValidRow()} and \texttt{C = dst.GetValidCol()}. Let \texttt{s\_j} be the per-column scalar taken from \texttt{src1} (one value per column).

For \texttt{0 \textless{}= i \textless{} R} and \texttt{0 \textless{}= j \textless{} C}:

\[
\mathrm{dst}_{i,j} = \min(\mathrm{src0}_{i,j}, s_j)
\]
}\par\smallskip
\paragraph{Notes.}
For FP types, NaN propagates: if either input is NaN the result is NaN.
\paragraph{Microcode site table.}
{\scriptsize
\begin{longtable}{@{}R{0.055\linewidth}L{0.23\linewidth}L{0.31\linewidth}L{0.22\linewidth}@{}}
\toprule
Seq & Micro instruction & WO[63:32] fields & WM[31:0] fields \\
\midrule
0 & \texttt{u\_vec.vlds} & \texttt{opc=0x17, x=0, fl=mem, seq=0, kind=b} & \texttt{entry=10, cnt=2, res=vec, var=0} \\
1 & \texttt{u\_vec.vlds} & \texttt{opc=0x17, x=0, fl=mem, seq=1, kind=b} & \texttt{entry=10, cnt=2, res=vec, var=0} \\
2 & \texttt{u\_vec.vmin} & \texttt{opc=0x1C, x=0, fl=alu, seq=2, kind=b} & \texttt{entry=10, cnt=2, res=vec, var=0} \\
3 & \texttt{u\_vec.vsts} & \texttt{opc=0x2A, x=0, fl=mem, seq=3, kind=b} & \texttt{entry=10, cnt=2, res=vec, var=0} \\
\bottomrule
\end{longtable}
}
\paragraph{Microcode body DSL.}
\begin{lstlisting}[style=ptocode]
def uop_tcolexpandmin(src0, src1, dst):
    dtype = dst.element_type
    valid_rows, valid_cols = dst.valid_shape

    for row in range(0, valid_rows, 1):
        remained = valid_cols
        for col in range(0, valid_cols, pto.get_lanes(dtype)):
            mask, remained = pto.make_mask(dtype, remained)
            lhs = pto.vlds(src0[row, col:])
            rhs = pto.vlds(src1[0, col:])
            result = pto.vmin(lhs, rhs, mask)
            pto.vsts(result, dst[row, col:], mask)
    return
\end{lstlisting}
\Needspace{0.34\textheight}
